%==============================================================================
%  Doing a master's thesis in ML / NLP -- a practical guide
%  Konstantinos Skianis, University of Ioannina
%  Compile:  latexmk -cd -pdf masters-thesis-ml-nlp-guide.tex
%==============================================================================
\documentclass[11pt,a4paper]{article}
\usepackage{guide-style}
\guidemeta{Master's thesis in ML / NLP}{Doing a Master's Thesis in Machine Learning or NLP}{Master's thesis supervision, ML/NLP research practice}

\begin{document}

%==============================================================================
\begin{titlepage}
\centering
\vspace*{2.0cm}
{\color{brandlight}\rule{\textwidth}{2pt}}\\[0.8cm]
{\Huge\bfseries\color{brand} Doing a Master's Thesis\\[0.25cm] in Machine Learning or NLP\par}
\vspace{0.7cm}
{\LARGE\color{ink} A Practical Guide\par}
\vspace{0.5cm}
{\large\color{mutegray} Finding a contribution, defending it, and what it becomes next\par}
\vspace{0.6cm}
{\color{brandlight}\rule{\textwidth}{2pt}}\\[1.1cm]

{\large\bfseries\color{ink} \guideauthor\par}
\vspace{0.18cm}
{\normalsize\color{mutegray} \guideaffiliation\par}
\vspace{0.25cm}
{\footnotesize\color{brandlight}\url{https://y3nk0.github.io}\par}
\vspace{1.1cm}

\begin{minipage}{0.85\textwidth}
\small\color{ink}
A master's thesis sits between an undergraduate project and a PhD chapter, and
it is judged by a different standard than either. You are expected to make a
\emph{contribution} --- something the reader did not know before --- but a
bounded one, finished inside two semesters. This guide is about choosing that
contribution, running experiments that survive scrutiny, writing it up, and
turning it into a paper, a PhD application, or a job.
\end{minipage}

\vfill
{\color{mutegray} October 2026}\\[0.3cm]
{\footnotesize\color{mutegray} Your department's own regulations --- length,
format, deadlines, ECTS, grading --- always take precedence over anything here.}
\end{titlepage}

%==============================================================================
\tableofcontents
\thispagestyle{fancy}
\clearpage

%==============================================================================
\section{What a master's thesis has to be}

Three projects, three standards. It helps to see them side by side, because
students routinely aim at the wrong one.

\begin{table}[h]
\centering
\small
\begin{tabularx}{\textwidth}{L{2.6cm} X X}
\toprule
 & \thead{Undergraduate} & \thead{Master's} \\
\midrule
\rowcolor{rowtint}
Question & Given, or chosen from a menu & \textbf{Yours}, argued for and scoped
with your supervisor \\
Bar & Done correctly, understood completely & Correct \emph{and} tells the
field something it did not know \\
\rowcolor{rowtint}
Related work & A summary of what exists & A \textbf{positioned} chapter: a
taxonomy, and where your work sits in it \\
Evidence & Baselines, a fair comparison, error analysis & All of that, plus
ablations, several seeds, and significance where it is claimed \\
\rowcolor{rowtint}
Typical output & A working system and a report & A report that could become a
\textbf{workshop paper} with a month's extra work \\
Length & 30--50 pages & 40--80 pages, by department rule \\
\bottomrule
\end{tabularx}
\end{table}

\begin{cbox}{The one-claim test}
Write the sentence your thesis exists to defend, in this shape:

\begin{quote}\small
\emph{For \textbf{[task/setting]}, \textbf{[what you did]} achieves
\textbf{[effect]} compared with \textbf{[baseline]}, because
\textbf{[mechanism]} --- shown by \textbf{[the experiment]}.}
\end{quote}

If you cannot fill every slot by the end of month two, you have a project, not
a thesis. The missing slot tells you what to work on next: usually
\emph{because} (you have numbers but no explanation) or \emph{compared with}
(you never built the honest baseline).
\end{cbox}

\begin{wbox}{What a PhD requires that you do not}
Novelty at the level of a main-conference paper, a programme of work spanning
years, and a survey of a whole subfield. Aiming there will cost you the
deadline. A master's contribution can be: a method applied somewhere it has
not been and evaluated properly; a careful reproduction that finds the limits
of a published claim; an ablation that identifies which component actually
matters; a new resource or benchmark for an under-served language or domain.
Each of those is publishable at a workshop and finishable in a semester.
\end{wbox}

%==============================================================================
\section{Finding a contribution}
\label{sec:contribution}

You need a question that is \emph{open}, \emph{answerable with your compute},
and \emph{interesting whichever way the answer falls}. That last property is
what protects you: if your idea works you have a result, and if it fails you
have a finding, provided the experiment was fair.

\subsection{Five routes that work}

\begin{description}
  \item[1. Extend one paper.] The most reliable route. Take a recent method
        you admire and push exactly one dimension: a new language or domain, a
        harder distribution shift, a tenth of the compute, a tenth of the
        labels, longer inputs, a different evaluation. The baseline is the
        paper; the contribution is where it breaks and why.
  \item[2. Reproduce and interrogate.] Rebuild a published result, then ask
        the question the authors did not: does it hold across seeds, datasets,
        model scales, prompts? Reproducibility work is publishable in its own
        right and teaches you more about experimental rigour than any course.
  \item[3. A shared task, with an angle.] SemEval and the workshop tasks at
        \mbox{*ACL} venues hand you data, a metric and a leaderboard
        (\S\ref{sec:tools} has the links). A master's thesis does not stop at
        the ranking: the thesis is the analysis of \emph{why} systems place
        where they do, which the system paper has no room for.
  \item[4. An applied or industry thesis.] A real problem from a company, a
        lab in another department, or a public body. The research question has
        to survive contact with the deployment constraint ---
        see \S\ref{sec:industry} before you sign anything.
  \item[5. Analysis or theory.] Probing what a representation encodes, a
        robustness study, an efficiency comparison measured honestly, a
        critique of a metric with a small human evaluation, or a bounded
        theoretical result. Keep the claim small and complete.
\end{description}

\subsection{Where the open questions are written down}

\begin{cbox}{Six places to look, in order of yield}
\begin{enumerate}
  \item \textbf{Limitations sections.} Mandatory at ACL venues: authors list
        what their method cannot do. It is a menu of theses.
  \item \textbf{Task overview papers.} The closing section says exactly what
        nobody solved this year.
  \item \textbf{Survey taxonomies.} Find the most recent survey and look for
        the boxes with one or two entries.
  \item \textbf{Workshop calls for papers.} This year's open problems, in
        bullet form, written by the people who will review you.
  \item \textbf{Leaderboards that have plateaued}, or where a simple baseline
        sits suspiciously close to the top --- a sign the task, metric or data
        deserves a second look.
  \item \textbf{Your supervisor's current grant.} Unglamorous and effective:
        the open questions are already funded, the data already exists, and
        the supervision will be attentive because the work matters to them.
\end{enumerate}
\end{cbox}

\subsection{The two-page proposal}

Before you write code, write two pages and get them signed off. This is longer
than the undergraduate version because the risks are larger.

\begin{center}
\small
\begin{tabularx}{0.97\textwidth}{L{3.6cm} X}
\toprule
\thead{Section} & \thead{What goes in it} \\
\midrule
\rowcolor{rowtint} Question & One sentence, open, answerable \\
Claim & The one-claim sentence above, with the slots filled as far as you can \\
\rowcolor{rowtint} Why now & Two or three references showing the gap is real
and current \\
Data & Name, size, licence, access route, and whether you already have it \\
\rowcolor{rowtint} Baselines & The trivial one, the published one, and where
its numbers come from \\
Method & What you will try, in three bullets, with the one knob you will vary \\
\rowcolor{rowtint} Evaluation & Metric, splits, how many seeds, and the test
you will use to claim a difference \\
Compute budget & GPU hours you can actually get, and the largest model that
fits \\
\rowcolor{rowtint} Milestones & Five dates: data, baseline, first result,
freeze, full draft \\
Risks & The three most likely ways this stalls \\
\rowcolor{rowtint} Fallback & The thesis you still have if the main idea fails
--- usually the comparison and the analysis \\
\bottomrule
\end{tabularx}
\end{center}

\begin{wbox}{Scope, the thing that actually sinks theses}
Not difficulty --- scope. Three questions, honestly: can you name the dataset
and metric today? Can you run one baseline end to end within three weeks? If
the main idea fails completely, is there still a thesis? Any ``no'' means
narrow it now, while narrowing is free.
\end{wbox}

%==============================================================================
\section{Industry, external and NDA-bound theses}
\label{sec:industry}

A thesis done inside a company can be excellent --- real data, real
constraints, a mentor who cares. It can also become unexaminable. Settle these
\emph{in writing, before you start}, with your academic supervisor in the
conversation.

\begin{cbox}{The six questions to ask on day one}
\begin{enumerate}
  \item \textbf{Can the thesis be published and defended openly?} A thesis is
        an examined document; most universities require it to be readable by
        examiners and often archived publicly.
  \item \textbf{Is there an embargo?} Many departments allow a delay of six or
        twelve months. Find out the rule and whether the company needs it.
  \item \textbf{What may appear?} Method, yes; raw data, often no; absolute
        numbers, sometimes no. Relative results (``+3.1 F1 over the internal
        baseline'') are usually negotiable when absolutes are not.
  \item \textbf{Who owns the code and the IP?} Say explicitly what you may put
        on GitHub. Agree it before you write it.
  \item \textbf{Who supervises what?} You need an academic supervisor
        responsible for the research standard and a company mentor for the
        problem. When they disagree, the degree requirements win.
  \item \textbf{What is the fallback if access is withdrawn?} Reorganisations
        and policy changes happen mid-thesis. A public dataset you can fall
        back to turns a catastrophe into a chapter.
\end{enumerate}
\end{cbox}

\begin{wbox}{Writing around confidential data}
If the data cannot be described, the thesis has to stand on the method and the
protocol: describe the data by its \emph{statistics and structure} (size,
label distribution, domain, annotation process) rather than its content, and
reproduce the core experiment on a public dataset so an examiner can check
your reasoning against something they can see. Do this from the start --- it
is painful to retrofit in month eight.
\end{wbox}

Personal data brings its own duties: GDPR applies to scraped social media,
user logs, recordings and surveys. Ask about ethics approval early; in some
departments it is required and takes weeks.

%==============================================================================
\section{Related work, the chapter that separates degrees}

The undergraduate version summarises papers. The master's version
\emph{positions} your work. Organise by theme, not by paper, and end by saying
what nobody has done --- which is, conveniently, what you are about to do.

\begin{cbox}{A bounded search you can finish in three weeks}
\begin{enumerate}
  \item \textbf{Seed set:} 5--10 anchor papers from your supervisor. Read them
        properly; their vocabulary becomes your queries.
  \item \textbf{Search systematically} across Google Scholar, Semantic
        Scholar, the ACL Anthology, OpenReview and DBLP --- and
        \textbf{log every query, date and hit count} in a spreadsheet. That
        log is what makes the chapter defensible rather than anecdotal.
  \item \textbf{Snowball both ways:} backward through reference lists,
        forward through ``cited by'' and \emph{search within citing articles}.
        Forward snowballing is how you find the last six months.
  \item \textbf{Screen in two passes:} title and abstract, then full text.
        Expect to screen 200--400 and keep 60--100.
  \item \textbf{Extract into a table}, one row per paper: key, venue, problem,
        method family, data, metric, headline result, compute, code, and
        whether it reproduces.
  \item \textbf{Cluster the rows} until the groups stop changing. That is your
        taxonomy, and it becomes one figure and the chapter's skeleton.
\end{enumerate}
\end{cbox}

Read in three passes (Keshav): five minutes to decide, an hour for figures and
method, several hours only for the five papers your thesis depends on.

%==============================================================================
\section{Tools, set up in week one}
\label{sec:tools}

The same stack as every other guide in this series; set it up once and stop
thinking about it.

\begin{itemize}
  \item \textbf{Zotero} (\url{https://www.zotero.org}) with the browser
        connector and \textbf{Better BibTeX}
        (\url{https://retorque.re/zotero-better-bibtex/}): pinned citation
        keys and auto-export into your thesis \code{refs.bib}.
        The quickest way to fill it: enable Scholar's BibTeX links
        (\emph{Settings} $\rightarrow$ \emph{Bibliography manager}), then
        \textbf{Cite} $\rightarrow$ \textbf{BibTeX}. Check every entry --- a
        bibliography full of ``arXiv preprint'' for published papers is the
        clearest sign nobody looked. Cleaner sources: the \textbf{ACL
        Anthology} (\url{https://aclanthology.org}), \textbf{DBLP}
        (\url{https://dblp.org}), or any DOI from the terminal with
        \code{curl -LH "Accept: application/x-bibtex" https://doi.org/<DOI>}.
  \item \textbf{LaTeX}, in \textbf{Overleaf}
        (\url{https://www.overleaf.com}) for supervisor comments, and
        \code{git} + \code{latexmk} locally. One sentence per line makes diffs
        readable; \code{latexdiff} shows your supervisor exactly what changed
        between drafts. \emph{Prism shuts down on 1 January 2027 --- do not
        start a thesis there, and export anything you have.}
  \item \textbf{One git repository} for code, configs, notes and the thesis.
        Commit daily.
  \item \textbf{Compute:} your lab's cluster or Slurm queue, department GPUs,
        or Colab and Kaggle notebooks for smaller runs. Ask in week one ---
        accounts take time, and the largest model you can train is a design
        constraint, not a detail.
  \item \textbf{Experiment tracking} from the first run (Weights \& Biases,
        MLflow, TensorBoard): config, metrics and the git commit, every time.
  \item \textbf{Shared tasks and deadlines:}
        \url{https://semeval.github.io},
        \url{https://paperswithcode.co/ai-deadlines} and
        \url{http://aideadlines.org/?sub=ML,CV,NLP}.
\end{itemize}

%==============================================================================
\section{Experiments that survive an examiner}

\subsection{The order that works}

\begin{enumerate}
  \item Get the data and read fifty examples by hand. Count labels, find
        duplicates, look for leakage between splits.
  \item Build the \textbf{trivial baseline} --- majority class, TF-IDF plus
        logistic regression, a rule --- and measure it. It anchors everything.
  \item Reproduce the \textbf{published baseline}. If you cannot match it,
        nothing downstream is interpretable, and finding out why is itself a
        result.
  \item Only then your method, one change at a time, each measured.
  \item \textbf{Ablate}: remove each component and show what it was worth.
        Start early; ablations always take longer than the headline run.
  \item \textbf{Error analysis}: categorise 50--100 mistakes by hand. This is
        where the \emph{because} in your one-claim sentence comes from.
\end{enumerate}

\subsection{Claiming a difference honestly}

\begin{cbox}{Seeds, spread and significance}
\begin{itemize}
  \item Run every configuration with at least \textbf{three seeds} --- five if
        the budget allows --- and report \textbf{mean $\pm$ standard
        deviation}. A single run that improves by 0.4 is noise, and an
        examiner will say so.
  \item When you claim a system is better, test it: a \textbf{paired bootstrap}
        or \textbf{permutation test} on the test predictions is standard in
        NLP, easy to implement, and makes no distributional assumptions.
  \item Comparing many systems on the same data inflates false positives ---
        say so, or correct for it.
  \item \textbf{Report compute}: GPU type, hours, parameters, and the cost of
        the hyperparameter search. A method that wins only after a hundred
        trials against a baseline tuned for one is not winning.
  \item Fix the splits \emph{before} you look at results, and touch the test
        set once.
\end{itemize}
\end{cbox}

\begin{wbox}{Negative results, handled properly}
``I tried X, it did not beat the baseline, here is the evidence and here is
why'' is a legitimate master's thesis --- provided the comparison was fair,
the implementation was checked against a reference, and the analysis explains
the failure. Many strong theses are exactly this. What is not acceptable is a
negative result from a bug you never looked for: before concluding, verify
your implementation reproduces a known number on a known dataset.
\end{wbox}

\subsection{The mistakes that cost a grade}

\begin{wbox}{Seen every year}
\begin{itemize}
  \item \textbf{Leakage.} Tuning on test, duplicates shared between splits, or
        preprocessing and feature selection fitted on the full dataset before
        splitting. It invalidates everything downstream.
  \item \textbf{An unfair comparison.} Your method tuned for a week, the
        baseline run on defaults for an hour. Examiners spot it immediately,
        and it is the fastest way to lose their trust in the rest.
  \item \textbf{A difference claimed from one seed.}
  \item \textbf{Accuracy on imbalanced data.} Report macro-F1, per-class
        results, or a confusion matrix.
  \item \textbf{The method chapter that cannot be reimplemented} because the
        details live only in the code.
  \item \textbf{A related-work chapter that is a list.} Paragraph per paper,
        no taxonomy, no positioning --- the clearest marker of an
        undergraduate thesis wearing a master's cover.
  \item \textbf{Writing started in the last month.} It shows, and it is the
        single strongest predictor of a mediocre grade.
  \item \textbf{Results that cannot be regenerated} because the configuration
        lived in your shell history.
\end{itemize}
\end{wbox}

\subsection{Keep a results table from day one}

One CSV or spreadsheet: date, commit, config, seed, dev score, test score
(only when you finally run it), and a one-line note on what changed and why.
Every table and figure in the thesis should regenerate from logged results
with one script. In month eight, ``which run produced Table 4?'' must have an
answer.

%==============================================================================
\section{Using an AI assistant}

An assistant removes mechanical work and can quietly remove your degree. Both
are true.

\begin{itemize}
  \item \textbf{Good uses:} explaining unfamiliar notation, getting someone
        else's repository to run, boilerplate and plotting code, refactoring,
        LaTeX debugging, tightening prose you wrote, and acting as a hostile
        examiner on your draft.
  \item \textbf{Set it up} in the thesis repository with a short memory file
        of your conventions, connected to your files and your Zotero library,
        and keep a dated \code{notes/} log it can read.
\end{itemize}

\begin{wbox}{Five rules}
\begin{enumerate}
  \item \textbf{You must be able to explain every line you submit.} The
        defence tests precisely this.
  \item \textbf{Never cite a paper you have not opened.} Models invent
        plausible references; fabricated citations end theses.
  \item \textbf{Never let it write your results or your analysis.} Numbers
        come from runs you did; interpretation is the degree.
  \item \textbf{Follow your university's policy} and declare assistance where
        a declaration is required. Ask your supervisor; do not guess.
  \item \textbf{Nothing confidential into a cloud tool} --- company data under
        NDA, personal data, or someone else's unpublished manuscript.
\end{enumerate}
\end{wbox}

%==============================================================================
\section{Writing it}

Start in month one. Writing is how you discover the hole in the argument, and
holes are cheap to fix early and ruinous to find in the last fortnight.

\begin{center}
\small
\begin{tabularx}{\textwidth}{L{3.3cm} X L{1.9cm}}
\toprule
\thead{Chapter} & \thead{What goes in it} & \thead{Pages} \\
\midrule
\rowcolor{rowtint}
Abstract & Problem, what you did, headline result, in one paragraph & 1 \\
Introduction & Context, the question, your contributions as a numbered list,
thesis outline & 4--6 \\
\rowcolor{rowtint}
Background & Only the concepts your chapters actually use & 6--12 \\
Related work & By theme, with a taxonomy figure; ends by positioning your work
& 8--12 \\
\rowcolor{rowtint}
Method & Precise enough to reimplement from the text alone & 8--12 \\
Experiments & Data, splits, metrics, baselines, protocol, then results &
10--18 \\
\rowcolor{rowtint}
Analysis & Ablations, error analysis, significance, what surprised you & 6--10 \\
Conclusion & Findings, limitations, future work & 2--4 \\
\rowcolor{rowtint}
References, appendix & Bibliography; hyperparameters, prompts, extra tables &
--- \\
\bottomrule
\end{tabularx}
\end{center}

\emph{Sizes are a sane default for a 40--80 page thesis; your department's
rules win.}

\begin{itemize}
  \item \textbf{Write the figures and tables first}, then the text explaining
        them. A table that needs three paragraphs of excuses means the
        experiment is wrong, not the prose.
  \item \textbf{Every chapter} opens with a sentence saying what it does and
        closes with one saying what it showed.
  \item \textbf{Contributions as a numbered list} in the introduction, each
        one checkable against a chapter.
  \item \textbf{Captions are self-contained}: a reader who looks only at your
        figures should still follow the thesis.
  \item \textbf{Cite properly, quote almost never.} Copying sentences --- from
        papers, from the web, or from a model --- without quotation marks and
        a citation is plagiarism, and it is checked.
  \item \textbf{Limitations, written honestly}, are a sign of maturity.
        Examiners find them anyway; better that you found them first.
  \item \textbf{A complete draft to your supervisor three weeks before the
        deadline.} Feedback on a partial draft is feedback on the wrong
        problem.
\end{itemize}

%==============================================================================
\section{Turning the thesis into a paper}

This is the step that makes a master's thesis count twice --- as a degree and
as a publication --- and most students skip it because they are exhausted.
Talk to your supervisor about it in month five, not after the defence.

\begin{itemize}
  \item \textbf{What to cut:} the background chapter, most of the related
        work, the narrative of what you tried. A paper keeps the claim, the
        method, the key experiments and the analysis --- eight pages, not
        sixty.
  \item \textbf{Where to send it:} a workshop at an \mbox{*ACL} or ML venue is
        the realistic first target; \textbf{Findings} and the \textbf{Student
        Research Workshop} are real, citable venues; a shared-task
        \textbf{system paper} if you took that route. The PhD guide in this
        series covers the venue landscape and the ACL Rolling Review.
  \item \textbf{Authorship:} agree the order early, in conversation, not by
        assumption. Typically you are first author and your supervisor last.
  \item \textbf{Timing:} the thesis deadline and the venue deadline rarely
        align. Submit the paper after you submit the thesis, using the
        examiners' comments as free reviews.
  \item \textbf{Release the artefacts:} code with a licence and a README that
        reproduces your main table, checkpoints on Hugging Face, and a
        \code{CITATION.cff} so people can cite it.
\end{itemize}

%==============================================================================
\section{The defence}

Usually 20--30 minutes plus questions, before a committee that has read the
thesis unevenly.

\begin{itemize}
  \item \textbf{Structure:} problem (2 slides) $\rightarrow$ why it is hard
        and what exists (2--3) $\rightarrow$ your claim (1) $\rightarrow$
        method (3--5) $\rightarrow$ results (4--5) $\rightarrow$ analysis and
        limitations (2--3) $\rightarrow$ conclusion (1). One idea per slide.
  \item \textbf{Rehearse out loud, timed, twice}, at least once to a person.
  \item \textbf{Know your numbers:} dataset sizes, your result, the baseline,
        the spread across seeds, training cost.
  \item \textbf{Expect:} why this baseline? is the improvement significant?
        what happens under distribution shift? why does it fail on class X?
        what would you do with six more months? what did you find hardest?
  \item \textbf{``I do not know''} is a fine answer when followed by ``here is
        how I would find out''. Guessing is not.
\end{itemize}

%==============================================================================
\section{Working with your supervisor}

At master's level the division of labour shifts: your supervisor sets the
direction and the standard, and you own the execution. The students who get
the most from supervision are not the ones who need the least help --- they
are the ones who ask precise questions.

\begin{cbox}{Decide yourself, bring to the meeting, escalate immediately}
\textbf{Decide yourself:} library and framework choices, code structure,
plotting, which of two equivalent baselines to run first, how to organise your
notes. Bringing these to a meeting wastes the half hour you have.\\[4pt]
\textbf{Bring to the meeting:} anything that changes the claim or the scope ---
a result that contradicts the plan, a baseline you cannot reproduce, a
dataset that turns out to be unusable, an idea that would add a month.\\[4pt]
\textbf{Escalate immediately, do not wait for the meeting:} losing data or
compute access, an ethics or NDA problem, a deadline you are going to miss,
and anything about your health or circumstances that affects the timeline.
Supervisors can nearly always help with these \emph{if told early} and almost
never at the end.
\end{cbox}

\begin{itemize}
  \item \textbf{Send the agenda the evening before} --- five bullets, as
        below. A supervisor who has seen your numbers before the meeting gives
        you an hour of thinking instead of ten minutes of catching up.
  \item \textbf{Disagree in writing, with evidence.} ``I ran the experiment
        you suggested; here is why I think the other route is better'' is a
        good conversation. Silently doing something else is not.
  \item \textbf{Ask for the standard explicitly}: ask to see a thesis they
        considered excellent, and what separated it from an average one.
  \item \textbf{Say when you are stuck.} Two days stuck is normal; two weeks
        stuck and silent is a supervision failure that is half yours.
  \item \textbf{What a supervisor cannot do:} write it for you, guarantee a
        result, or rescue a thesis in the final fortnight. Plan accordingly.
\end{itemize}

%==============================================================================
\section{A timetable}

A nine-month plan; compress to six by halving months 2--4, and put every
milestone in a calendar.

\begin{center}
\small
\begin{tabularx}{\textwidth}{L{2.0cm} X}
\toprule
\thead{Month} & \thead{Done by the end of it} \\
\midrule
\rowcolor{rowtint}
\textbf{1} & Topic agreed; two-page proposal signed off; tools installed; data
in hand and inspected; 15--20 papers read; trivial baseline running. \\
\textbf{2} & Published baseline reproduced; experiment tracking in place;
search log running; related-work taxonomy drafted. \\
\rowcolor{rowtint}
\textbf{3--4} & Method implemented; first real results; method chapter
drafted; one-claim sentence filled in completely. \\
\textbf{5} & Ablations and seeds; significance tests; publication plan agreed
with your supervisor; experiments chapter drafted. \\
\rowcolor{rowtint}
\textbf{6--7} & Error analysis; analysis chapter; results frozen; all tables
regenerated by script. \\
\textbf{8} & \textbf{Complete draft} to your supervisor; artefacts cleaned and
released. \\
\rowcolor{rowtint}
\textbf{9} & Revisions; every citation checked; defence slides; two
rehearsals; submission. \\
\bottomrule
\end{tabularx}
\end{center}

\begin{cbox}{Weekly rhythm}
Meet weekly, and send five bullets the evening before: \textbf{what I did}
$\cdot$ \textbf{what I found} $\cdot$ \textbf{what is blocked} $\cdot$
\textbf{what I propose next} $\cdot$ \textbf{what I need from you}. Keep a
shared running document. ``I am stuck'' sent on Tuesday is worth three weeks
of silence.
\end{cbox}

%==============================================================================
\section{What it becomes next}

\begin{description}
  \item[A PhD application.] Your thesis is the main evidence that you can do
        research. Give your referees the abstract, the results and your plans;
        ask while the work is fresh. A workshop paper from the thesis moves
        you from ``promising'' to ``has published''.
  \item[A job.] The repository is the portfolio: a README that reproduces your
        main table, honest limitations, and clean code say more than any
        bullet point on a CV. Expect to be asked, in interviews, exactly what
        your examiners asked.
  \item[Both.] They are not exclusive, and an industry year before a PhD is
        a perfectly good path.
\end{description}

%==============================================================================
\section{Checklists}

\subsection*{Month 1}
\begin{multicols}{2}
\begin{itemize}
  \item Supervisor meeting slot fixed; agenda format agreed
  \item Two-page proposal signed off
  \item Zotero + Better BibTeX; auto-export into \code{refs.bib}
  \item Git repository with \code{notes/} and a README
  \item Data obtained; licence and ethics checked
  \item Compute access working; a real job submitted
  \item Trivial baseline measured
  \item If industry: the six questions answered \emph{in writing}
\end{itemize}
\end{multicols}

\subsection*{Before freezing results}
\begin{itemize}
  \item Trivial baseline, published baseline and your method on identical
        splits
  \item Three or more seeds; mean $\pm$ std reported everywhere
  \item Significance test for every comparison you claim
  \item Compute and hyperparameter search reported for all systems
  \item Test set used once; no tuning on it
  \item Ablations complete; errors categorised by hand
  \item Every table and figure regenerated by script from logged results
  \item Code runs from a clean clone, following your own README
\end{itemize}

\subsection*{Before submitting}
\begin{itemize}
  \item One-claim sentence appears, in plain words, in the abstract and the
        conclusion
  \item Contributions list in the introduction, each checkable against a
        chapter
  \item Every citation checked against the real paper --- no invented
        references
  \item Limitations written honestly
  \item Department format, length and declaration requirements met
  \item Complete draft seen by your supervisor at least three weeks earlier
  \item Slides rehearsed twice, timed; demo recorded as a fallback
  \item Publication plan agreed: venue, authorship, timing
\end{itemize}

%==============================================================================
\appendix
\section{Link directory}

\begin{longtable}{L{4.3cm} L{10.5cm}}
\toprule
\thead{What} & \thead{Where} \\
\midrule
\endfirsthead
\toprule
\thead{What} & \thead{Where} \\
\midrule
\endhead
\multicolumn{2}{l}{\textbf{\color{brand}Literature}}\\
Google Scholar & \url{https://scholar.google.com} \\
Semantic Scholar & \url{https://www.semanticscholar.org} \\
ACL Anthology & \url{https://aclanthology.org} \\
DBLP & \url{https://dblp.org} \\
OpenReview & \url{https://openreview.net} \\
Connected Papers & \url{https://www.connectedpapers.com} \\
Papers with Code & \url{https://paperswithcode.co} \\
\midrule
\multicolumn{2}{l}{\textbf{\color{brand}Tasks and deadlines}}\\
SemEval & \url{https://semeval.github.io} \\
Conference deadlines & \url{https://paperswithcode.co/ai-deadlines} \\
 & \url{http://aideadlines.org/?sub=ML,CV,NLP} \\
\midrule
\multicolumn{2}{l}{\textbf{\color{brand}Tools}}\\
Zotero & \url{https://www.zotero.org} \\
Better BibTeX & \url{https://retorque.re/zotero-better-bibtex/} \\
Overleaf & \url{https://www.overleaf.com} \\
Hugging Face & \url{https://huggingface.co} \\
Weights \& Biases & \url{https://wandb.ai} \\
\midrule
\multicolumn{2}{l}{\textbf{\color{brand}The companion guides}}\\
Undergraduate thesis, PhD & \url{https://y3nk0.github.io/guides/} \\
\bottomrule
\end{longtable}

\section{The two-page proposal, as a form}

\begin{quote}\small
\textbf{Title.} Working title.\\[3pt]
\textbf{Question.} The one sentence this thesis answers.\\[3pt]
\textbf{Claim.} \emph{For \rule{2.2cm}{0.4pt}, \rule{2.2cm}{0.4pt} achieves
\rule{2cm}{0.4pt} compared with \rule{2cm}{0.4pt}, because
\rule{2cm}{0.4pt} --- shown by \rule{2.4cm}{0.4pt}.}\\[3pt]
\textbf{Why now.} Two or three references showing the gap.\\[3pt]
\textbf{Route.} Extend a paper / reproduce / shared task / applied /
analysis (\S\ref{sec:contribution}).\\[3pt]
\textbf{Data.} Name, size, licence, access, already in hand?\\[3pt]
\textbf{Baselines.} Trivial; published; source of its numbers.\\[3pt]
\textbf{Method.} Three bullets, and the one knob you vary.\\[3pt]
\textbf{Evaluation.} Metric, splits, seeds, significance test.\\[3pt]
\textbf{Compute.} Hardware, hours, largest model that fits.\\[3pt]
\textbf{Milestones.} Data \rule{1.4cm}{0.4pt}; baseline
\rule{1.4cm}{0.4pt}; first result \rule{1.4cm}{0.4pt}; freeze
\rule{1.4cm}{0.4pt}; draft \rule{1.4cm}{0.4pt}.\\[3pt]
\textbf{Risks.} Three ways this stalls.\\[3pt]
\textbf{Fallback.} The thesis you still have if the idea fails.\\[3pt]
\textbf{If industry.} Publication, embargo, what may appear, IP, supervision,
fallback access.
\end{quote}

\vspace{0.5cm}
\begin{cbox}{The shortest version of all of this}
Write the one-claim sentence in month one and let it tell you what is missing.
Reproduce the baseline before you innovate. Change one thing at a time, with
enough seeds to know it is real. Write from the first month, release what you
built, and plan the paper before you are exhausted. A bounded contribution,
finished and defended, beats an ambitious one abandoned in month eight.
\end{cbox}

\end{document}
